\chapter*{Dos propiedades que conectan información y modelado}
\addcontentsline{toc}{chapter}{Propiedades que conectan información y modelado}

\section*{La información no aumenta al agrupar categorías}
\addcontentsline{toc}{section}{Agregación y divergencia KL}
Supongamos que un experimento observa tres categorías, pero un instrumento
solo informa si ocurrió $a$ o si ocurrió alguna de las otras dos.
La transformación $f: \{a,b,c\}\to\{a,\mathrm{otro}\}$ agrupa
los resultados $b$ y $c$. Para
$p=(1/2,1/4,1/4)$ y $q=(1/4,1/2,1/4)$, teníamos
$D_{\rm KL}(p\|q)=1/4$ bit. Después de la agregación,
\[
 p_f=(1/2,1/2),\qquad q_f=(1/4,3/4),
\]
y la divergencia entre las distribuciones observables es
\[
\begin{split}
D_{\rm KL}(p_f\|q_f)
&=\frac12\log_2\frac{1/2}{1/4}
+\frac12\log_2\frac{1/2}{3/4}\\
&=\frac12+\frac12\log_2\frac23
\approx0.207519\ \text{bits}.
\end{split}
\]
Ha disminuido respecto del $0.25$ original. Este ejemplo particular
ilustra una propiedad general de la KL denominada \emph{procesamiento de
datos}: al aplicar una transformación determinista común a los
resultados, no aumenta la divergencia.

\begin{figure}[H]\centering
\begin{tikzpicture}[>=Stealth,font=\small,
 node/.style={draw=azul,rounded corners,fill=azulclaro,minimum width=17mm,minimum height=8mm}]
\node[node](a)at(0,1){$a$};\node[node](b)at(0,0){$b$};\node[node](c)at(0,-1){$c$};
\node[node](yes)at(6,0.8){$a$};\node[node](oth)at(6,-0.8){otro};
\draw[->] (a)--(yes);\draw[->] (b)--(oth);\draw[->] (c)--(oth);
\node at (3,-1.8) {La observación final olvida la distinción entre $b$ y $c$.};
\end{tikzpicture}
\caption{La agregación pierde información sobre la diferencia entre dos leyes.}
\end{figure}

\textbf{Justificación.} Si $A_y=f^{-1}(y)$, la desigualdad de suma logarítmica
permite escribir
\[
\sum_{i\in A_y}p_i\ln\frac{p_i}{q_i}
\ge\left(\sum_{i\in A_y}p_i\right)
\ln\frac{\sum_{i\in A_y}p_i}{\sum_{i\in A_y}q_i}.
\]
Sumando en las categorías agregadas se obtiene
$D_{\rm KL}(p\|q)\ge D_{\rm KL}(p_f\|q_f)$.
La prueba exige las convenciones habituales de soporte, y no afirma
que ambas divergencias tengan que ser diferentes: podrían coincidir.
Este resultado es útil en investigación aplicada cuando un reporte
técnico se reduce a etiquetas de eventos; esa reducción puede eliminar
diferencias que existían a nivel textual.

\section*{Por qué ampliar el contexto puede reducir la incertidumbre óptima}
\addcontentsline{toc}{section}{Condicionar reduce la entropía en promedio}
Sea $X$ el próximo símbolo, $Y$ el último símbolo observado y $Z$
una porción anterior del contexto. La teoría de la información
establece
\[
H(X\mid Y,Z)\le H(X\mid Y).
\]
No es una afirmación acerca de que \emph{cualquier modelo entrenado}
con más contexto rendirá mejor. Es una desigualdad sobre las verdaderas
entropías condicionales de una distribución fija.

Para justificarla, definamos la información mutua condicional:
\[
I(X;Z\mid Y)
=\mathbb E_{Y,Z}\Bigl[D_{\rm KL}\bigl(
 p(X\mid Y,Z)\,\|\,p(X\mid Y)
\bigr)\Bigr].
\]
Por Gibbs, el integrando es no negativo en cada contexto del soporte.
Al expandir los logaritmos también se obtiene
\[
I(X;Z\mid Y)=H(X\mid Y)-H(X\mid Y,Z).
\]
Combinando ambas igualdades deducimos la desigualdad.
La igualdad ocurre si $X$ y $Z$ son condicionalmente independientes dado $Y$.

\begin{ejemplo}[Una dependencia que el último símbolo no explica]
En la construcción $X_1$ y $X_2$ Bernoulli$(1/2)$ independientes y
$X_3=X_1$, se tiene
\[
H(X_3\mid X_2)=1\ \text{bit},
\qquad H(X_3\mid X_1,X_2)=0.
\]
La historia más amplia elimina toda incertidumbre \emph{verdadera}
sobre $X_3$, mientras que el bigrama conserva un bit de incertidumbre.
En muestras pequeñas, sin embargo, estimar más contextos puede aumentar
la varianza y empeorar la predicción en prueba. Allí reaparece el
compromiso entre capacidad del modelo, datos y suavizado.
\end{ejemplo}
La transición de este resultado a los Transformers es conceptual:
el acceso a contexto amplio \emph{permite} representar dependencias
que los bigramas descartan, pero aprenderlas de forma fiable exige
datos, parametrización, optimización y evaluación, asuntos de S02--S05.
